% ==================== 附录 ====================
\newpage
\appendix
\ctexset{
  section = {
    name   = {附录},
    number = \Alph{section},
    format = \heiti\zihao{3}\raggedright,
  },
}

\section{需求—模块—图—界面 追溯矩阵}

本附录把实验一的 36 条需求一次性映射到模块、UML 图与界面，用于核对「有没有哪条需求在设计里落了空」。「主要图」列取自第 5 章的 UML 图清单（表~\ref{tab:uml_list}）；界面号按第 6 章的 01--16 屏；「—」表示该需求不在界面层体现（属于后端能力或范围约束）。

{\tabfont
\begin{xltabular}{\textwidth}{@{}L{1.75cm}L{2.75cm}L{2.65cm}L{3.75cm}X@{}}
\caption{需求追溯矩阵}\label{tab:app_trace}\\
\toprule
需求 & 名称 & 模块 & 主要图 & 界面 \\
\midrule
\endfirsthead
\toprule
需求 & 名称 & 模块 & 主要图 & 界面 \\
\midrule
\endhead
\midrule
\multicolumn{5}{r}{\zihao{6}续下页}\\
\endfoot
\bottomrule
\endlastfoot
REQ-001 & 角色定义与最小权限 & M1 & 用例图 $\times$4、包图 & — \\
REQ-002 & 顾客查询与预订 & M2/M4/M5 & 类图、活动图 1 & 01, 02, 03, 04 \\
REQ-003 & 退房与评价提交 & M7/M8 & 状态图·订单、用例图·顾客 & 06, 07, 10 \\
REQ-004 & 前台办理入住 & M7 & 用例图·前台、状态图·房间 & 09 \\
REQ-005 & 房间状态机 & M7 & 状态图·房间 & 08, 09, 10 \\
REQ-006 & 店长门店管理 & M2/M3/M4/M7 & 用例图·店长、类图 & 12, 14 \\
REQ-007 & 平台地区负责人管理 & M1/M3/M8 & 用例图·平台 & 15, 16 \\
REQ-008 & 订单状态机 & M5 & 状态图·订单 & 06 \\
REQ-009 & 取消与退款 & M5/M6 & 类图、时序图·超时释放 & 05, 06, 16 \\
REQ-010 & 重复提交防护 & M5 & 活动图 1、类图 & 04 \\
REQ-011 & 并发抢房 & M4 & 时序图·并发抢房、活动图 1 & 03 \\
REQ-012 & 改价与下单隔离 & M3 & 时序图·改价互斥 & 12, 13 \\
REQ-013 & 库存与可售上限 & M4 & 类图、组件图 & 02, 03, 12 \\
REQ-014 & 操作日志 & M9 & 组件图 & 13, 16 \\
REQ-015 & 性能 & 横切 & 部署图 & — \\
REQ-016 & 安全 & 横切 & 包图、部署图 & — \\
REQ-017 & 可靠性与数据一致 & 横切 & 部署图 & — \\
REQ-018 & 易用性 & 横切 & 用例图·前台 & 09 \\
REQ-019 & 范围约束 & 约束 & — & — \\
REQ-020 & 支付超时释放 & M5/M6 & 状态图·订单、时序图·超时释放 & 04, 05 \\
REQ-021 & 取消截止与前台代取消 & M5 & 状态图·订单、用例图·顾客 & 05, 06 \\
REQ-022 & 未到店标记与迟入住 & M5/M7 & 状态图·订单 & 08, 10 \\
REQ-023 & 价格快照与定价依据 & M3/M5 & 类图（明细单价快照） & 03, 04 \\
REQ-024 & 房号分配 & M7 & 类图（占房记录）、用例图·前台 & 09 \\
REQ-025 & 房态变更工单 & M7 & 类图（工单）、用例图·前台 & 10, 13 \\
REQ-026 & 员工信息可见范围 & M1 & 用例图 $\times$4 & 14, 15 \\
REQ-027 & 输入校验与错误提示 & 横切 & 活动图 1 & 04, 11 \\
REQ-028 & 支付与订单一致性 & M6/M5 & 时序图·支付回调、时序图·超时释放 & 16 \\
REQ-029 & 通知失败处理 & M10 & 组件图 & — \\
REQ-030 & 数据保留与删除审批 & M9 & 用例图·平台 & 16 \\
REQ-031 & 订单变更 & M3/M4/M5 & 用例图·顾客 & 06 \\
REQ-032 & 多房型多晚次组合 & M4/M5 & 类图（明细唯一约束） & 11 \\
REQ-033 & 线上可售上限与线下保留 & M2/M4 & 组件图、类图 & 11, 12, 13 \\
REQ-034 & 维修中与可售上限联动 & M2/M4/M7 & 状态图·房间、组件图 & 08, 12, 13 \\
REQ-035 & 订单排房时限 & —— & —— & 待确认，不覆盖 \\
REQ-036 & 账号注册与登录 & —— & —— & 待确认，不覆盖 \\
\end{xltabular}
}

从这张表可以读出三件事：第一，36 条需求中 34 条有设计落点，2 条因实验一未确认而明确不覆盖；第二，被引用最多的图是类图与组件图，说明「数据形状」与「模块边界」是本设计的两根主线；第三，有四条需求（REQ-015、016、017、019）在图与界面上都没有单独落点，它们是被全系统共享的非功能约束——这类需求不应硬塞给某个模块，写清「为什么不单独落点」比强行分配更有说服力。

\needspace{9cm}
\section{提交物清单}

提交包解开后是一层根目录加两个子目录。\texttt{UML/} 下每个 \texttt{.png} 都配有一份同名的 \texttt{.puml}（16 组、32 个文件），下表只列 \texttt{.png}。

{\zihao{5}\ttfamily\raggedright
\noindent 实验2\_241220031\_曹继贤/\par
\noindent\hspace*{2em}实验报告.pdf\par
\noindent\hspace*{2em}UML/\par
\noindent\hspace*{4em}UML类图.png\qquad UML包图.png\par
\noindent\hspace*{4em}UML组件图.png\qquad UML部署图.png\par
\noindent\hspace*{4em}UML用例图\_顾客端.png\qquad UML用例图\_前台端.png\par
\noindent\hspace*{4em}UML用例图\_店长端.png\qquad UML用例图\_平台端.png\par
\noindent\hspace*{4em}UML状态图\_订单.png\qquad UML状态图\_房间.png\par
\noindent\hspace*{4em}UML活动图\_下单与锁库存.png\qquad UML活动图\_支付回调与订单生效.png\par
\noindent\hspace*{4em}UML时序图\_并发抢房.png\qquad UML时序图\_改价互斥.png\par
\noindent\hspace*{4em}UML时序图\_支付回调.png\qquad UML时序图\_超时释放与对账.png\par
\noindent\hspace*{2em}UI界面/\par
\noindent\hspace*{4em}01\_home\_search.png\qquad 02\_search\_result.png\par
\noindent\hspace*{4em}03\_hotel\_detail.png\qquad 04\_order\_confirm.png\par
\noindent\hspace*{4em}05\_payment.png\qquad 06\_my\_orders.png\par
\noindent\hspace*{4em}07\_review\_submit.png\qquad 08\_room\_board.png\par
\noindent\hspace*{4em}09\_checkin.png\qquad 10\_housekeeping.png\par
\noindent\hspace*{4em}11\_walkin\_order.png\qquad 12\_price\_calendar.png\par
\noindent\hspace*{4em}13\_approval\_center.png\qquad 14\_roster\_review.png\par
\noindent\hspace*{4em}15\_region\_overview.png\qquad 16\_price\_range\_staffing.png\par
}

\section{图的绘制方式}

全部 16 张 UML 图由 PlantUML 绘制，源文件（\texttt{.puml}）与导出图（\texttt{.png}）同名放在 \texttt{UML/} 下，可用以下命令重新生成：

\begin{quote}
\zihao{5}\ttfamily
java -jar plantuml.jar -charset UTF-8 -tpng UML/*.puml
\end{quote}

统一的绘图约定：白底、深灰线、单一强调色；中文字体统一为 Microsoft YaHei；结构图与行为图宽度按内容自适应，时序图与活动图固定较窄的画幅以便插入正文。16 屏界面用 HTML 与 CSS 编写后按 2 倍像素密度导出 PNG（顾客端 780 $\times$ 1688，后台 2880 $\times$ 1800），因此插入报告后文字仍可辨认。
